% ============================================================
%  Baker N-Series Geometric Acceleration Framework
%  Assembled from collaborative research session
%  Author: Wayne Baker ("The Mystro")
% ============================================================

\documentclass[12pt]{article}
\usepackage{amsmath, amsthm, amssymb}
\usepackage{geometry}
\geometry{margin=1in}
\usepackage{hyperref}

% --- Theorem environments ---
\newtheorem{theorem}{Theorem}[section]
\newtheorem{lemma}[theorem]{Lemma}
\newtheorem{proposition}[theorem]{Proposition}
\newtheorem{corollary}[theorem]{Corollary}
\theoremstyle{remark}
\newtheorem{remark}{Remark}[section]

\title{\textbf{Geometric Acceleration of the Archimedean Limit:\\
Symmetry, Cancellation, and the N-Series Framework}}
\author{Wayne Baker}
\date{}

\begin{document}
\maketitle

% ============================================================
\begin{abstract}
We develop a systematic framework for accelerating convergence of the classical
Archimedean polygon sequence to $\pi$ using multi-scale linear combinations
whose coefficients preserve the exact limiting value. The even-power error
structure underlying the acceleration is shown to arise necessarily from
the odd symmetry of chord geometry under orientation reversal. A general
$r$-scale theorem is proved, yielding $O(N^{-2r-2})$ convergence at each
stage while preserving $\pi_g = \pi$ exactly. The same symmetry principle
is shown to govern the cubic refinement law in approximate angle trisection.
\end{abstract}

% ============================================================
\section{Introduction}

The Archimedean approximation
\[
  P_N = N\sin\!\left(\frac{\pi}{N}\right) \longrightarrow \pi
\]
converges at rate $O(N^{-2})$. The central observation of this paper is that the error expansion
\[
  P_N = \pi + \frac{C_2}{N^2} + \frac{C_4}{N^4} + \cdots
\]
contains \emph{only even powers} of $N^{-1}$, a fact determined by geometric symmetry rather than incidental algebraic structure. This structural property is derived rigorously in Section~\ref{sec:symmetry}, where it is shown to follow from the odd symmetry of the chord law under orientation reversal.

This symmetry-induced structure permits systematic cancellation: linear combinations of the sequence evaluated at multiple scales can be constructed to eliminate successive error terms while preserving the exact limit. The central thesis of this paper is that geometric symmetry determines the algebraic structure of approximation error, and that this structure admits exact, limit-preserving acceleration through multi-scale cancellation.

We call this family of operators the \emph{Baker N-series acceleration operators}. The paper proceeds in three stages: (i) the foundational symmetry principle, (ii) the general $r$-scale theorem with explicit Lagrange-form coefficients, and (iii) the connection to the cubic trisection refinement law.

% ============================================================
\section{Foundational Symmetry Principle}
\label{sec:symmetry}
\subsection{Orientation and Signed Geometric Response}

Geometric constructions involving angular displacement naturally depend on
orientation. If the orientation is reversed, the signed geometric response
changes sign. Thus, for a small angular parameter $\theta$, the fundamental
geometric response $f(\theta)$ satisfies
\[
  f(-\theta) = -f(\theta),
\]
and is therefore an odd function of $\theta$. A canonical example is the
chord of a circle:
\[
  c(\theta) = 2R\sin\!\left(\frac{\theta}{2}\right),
\]
which is odd in $\theta$.

\subsection{Normalization and Even Invariants}

While the raw geometric response is odd, many geometric approximations depend
on the normalized quantity $f(\theta)/\theta$. Since $f(\theta)$ is odd, the
quotient is even:
\[
  \frac{f(-\theta)}{-\theta} = \frac{f(\theta)}{\theta}.
\]
For the chord,
\[
  \frac{c(\theta)}{\theta} = R\,\frac{\sin(\theta/2)}{\theta/2},
\]
which is an even analytic function of $\theta$. Its Taylor expansion therefore
contains only even powers:
\[
  \frac{c(\theta)}{\theta}
  = R\!\left(1 - \frac{\theta^2}{24} + \frac{\theta^4}{1920}
    - \frac{\theta^6}{322560} + \cdots\right).
\]

\begin{lemma}[Odd-to-Even Conversion in Chord Geometry]
Let $f(x)$ be analytic near $x=0$ and odd, so that $f(-x)=-f(x)$.
Then the normalized function $g(x)=f(x)/x$ is analytic and even near $x=0$.
In particular, its Taylor expansion contains only even powers of $x$.
\end{lemma}

\begin{proof}
Since $f$ is odd, its Taylor expansion has the form
$f(x)=a_1x+a_3x^3+a_5x^5+\cdots$.
Dividing by $x$ gives $g(x)=a_1+a_3x^2+a_5x^4+\cdots$,
which contains only even powers. Hence $g$ is even.
\end{proof}

\subsection{Even-Power Structure of the Archimedean Sequence}

\begin{proposition}[Even-Power Expansion of $P_N$]
The Archimedean polygon sequence satisfies
\[
  P_N = N\sin\!\left(\frac{\pi}{N}\right)
      = \pi\,\frac{\sin(\pi/N)}{\pi/N},
\]
and hence
\[
  P_N = \pi - \frac{\pi^3}{6N^2} + \frac{\pi^5}{120N^4}
      - \frac{\pi^7}{5040N^6} + \cdots.
\]
In particular, the asymptotic expansion involves only even powers of $N^{-1}$.
\end{proposition}

\begin{proof}
Set $x = \pi/N$. Then $P_N = \pi\cdot(\sin x)/x$.
Since $\sin x$ is odd, the Lemma shows $(\sin x)/x$ is even, giving
\[
  \frac{\sin x}{x} = 1 - \frac{x^2}{6} + \frac{x^4}{120}
                    - \frac{x^6}{5040} + \cdots.
\]
Substituting $x=\pi/N$ yields the stated expansion.
\end{proof}

\begin{remark}
The even-power structure is forced by geometry: it is not an arithmetic
coincidence but the direct consequence of the odd symmetry of the chord law
under orientation reversal.
\end{remark}

% ============================================================
\section{Exactness of the Accelerated Geometric Limit}

\begin{theorem}[Limit Preservation under N-Series Acceleration]
\label{thm:two-scale}
Let $P_N = N\sin\!\left(\frac{\pi}{N}\right)$, so that $\lim_{N\to\infty}P_N = \pi$.
For a fixed integer $k > 1$, define the accelerated sequence
\[
  \widehat{P}_N^{(k)}
  = P_{kN} + \frac{P_{kN} - P_N}{k^2 - 1}.
\]
Then
\[
  \lim_{N \to \infty} \widehat{P}_N^{(k)} = \pi.
\]
Consequently, if $\pi_g$ is defined as the limit of any such N-series
accelerated geometric construction, then $\pi_g = \pi$.
\end{theorem}

\begin{proof}
Since $P_N \to \pi$, also $P_{kN} \to \pi$. Rewrite:
\[
  \widehat{P}_N^{(k)}
  = \frac{k^2}{k^2-1}P_{kN} - \frac{1}{k^2-1}P_N.
\]
Taking limits term by term,
\[
  \lim_{N\to\infty}\widehat{P}_N^{(k)}
  = \frac{k^2}{k^2-1}\pi - \frac{1}{k^2-1}\pi
  = \frac{k^2-1}{k^2-1}\,\pi = \pi. \qedhere
\]
\end{proof}

\begin{remark}
The decisive observation is that the coefficients $k^2/(k^2-1)$ and
$-1/(k^2-1)$ sum to $1$. This unity-sum condition is the sole requirement
for limit preservation; the rate of convergence is a separate matter
governed by the cancellation structure.
\end{remark}

The following general theorem makes the mechanism explicit.

\begin{theorem}[Linear Invariance of Convergent Sequences]
\label{thm:linear-invariance}
Let $\{a_N\}$ be a sequence with $\lim_{N\to\infty}a_N = L$.
Let $m_1,\dots,m_r$ be fixed positive integers and define
$A_N = \sum_{j=1}^r c_j\,a_{m_jN}$,
where $\sum_{j=1}^r c_j = 1$.
Then $\lim_{N\to\infty}A_N = L$.
\end{theorem}

\begin{proof}
Since $a_N \to L$, we have $a_{m_jN} \to L$ for each $j$.
Hence $\lim_{N\to\infty}A_N = \sum_j c_j L = \bigl(\sum_j c_j\bigr)L = L$.
\end{proof}

% ============================================================
\section{Multi-Scale Acceleration Theorems}

\subsection{Three-Scale Acceleration: $O(N^{-6})$}

\begin{theorem}[Three-Scale Limit-Preserving Acceleration]
\label{thm:three-scale}
Let
\[
  a_N = L + \frac{C_2}{N^2} + \frac{C_4}{N^4} + \frac{C_6}{N^6}
        + O(N^{-8}),
\]
and let $k, \ell > 1$ be distinct integers. Define
\[
  A_N = c_1 a_N + c_2 a_{kN} + c_3 a_{\ell N},
\]
where
\[
  c_1 = \frac{1}{k^2\ell^2},\qquad
  c_2 = -\frac{k^2}{(k^2-1)(\ell^2-k^2)},\qquad
  c_3 = \frac{\ell^2}{(\ell^2-1)(\ell^2-k^2)}.
\]
Then $A_N = L + O(N^{-6})$ and $\lim_{N\to\infty}A_N = L$.
\end{theorem}

\begin{proof}
Substitute the asymptotic expansion:
\[
  A_N = (c_1+c_2+c_3)L
      + \Bigl(c_1+\tfrac{c_2}{k^2}+\tfrac{c_3}{\ell^2}\Bigr)\frac{C_2}{N^2}
      + \Bigl(c_1+\tfrac{c_2}{k^4}+\tfrac{c_3}{\ell^4}\Bigr)\frac{C_4}{N^4}
      + O(N^{-6}).
\]
The stated coefficients satisfy
$c_1+c_2+c_3=1$,
$c_1+c_2/k^2+c_3/\ell^2=0$,
$c_1+c_2/k^4+c_3/\ell^4=0$.
Hence the $N^{-2}$ and $N^{-4}$ terms vanish, leaving $A_N = L+O(N^{-6})$.
\end{proof}

\begin{corollary}[Three-Scale Polygon Acceleration, $k=2$, $\ell=4$]
Let $P_N = N\sin(\pi/N)$. Then
\[
  \widetilde{P}_N
  = \tfrac{1}{64}P_N - \tfrac{4}{9}P_{2N} + \tfrac{64}{45}P_{4N}
  = \pi + O(N^{-6}),\qquad
  \lim_{N\to\infty}\widetilde{P}_N = \pi.
\]
\end{corollary}

\subsection{Four-Scale Acceleration: $O(N^{-8})$}

\begin{theorem}[Four-Scale Limit-Preserving Acceleration]
\label{thm:four-scale}
Let
\[
  a_N = L + \frac{C_2}{N^2} + \frac{C_4}{N^4}
        + \frac{C_6}{N^6} + \frac{C_8}{N^8} + O(N^{-10}).
\]
Define
\[
  A_N = c_0 a_N + c_1 a_{2N} + c_2 a_{4N} + c_3 a_{8N},
\]
where
\[
  c_0 = -\frac{1}{2835},\quad
  c_1 = \frac{4}{135},\quad
  c_2 = -\frac{64}{135},\quad
  c_3 = \frac{4096}{2835}.
\]
Then $A_N = L + O(N^{-8})$ and $\lim_{N\to\infty}A_N = L$.
\end{theorem}

\begin{proof}
Substituting the asymptotic expansion gives
\begin{align*}
  A_N &= (c_0+c_1+c_2+c_3)L \\
      &\quad+ \Bigl(c_0+\tfrac{c_1}{4}+\tfrac{c_2}{16}+\tfrac{c_3}{64}\Bigr)
             \frac{C_2}{N^2} \\
      &\quad+ \Bigl(c_0+\tfrac{c_1}{16}+\tfrac{c_2}{256}+\tfrac{c_3}{4096}\Bigr)
             \frac{C_4}{N^4} \\
      &\quad+ \Bigl(c_0+\tfrac{c_1}{64}+\tfrac{c_2}{4096}
             +\tfrac{c_3}{262144}\Bigr)\frac{C_6}{N^6}
             + O(N^{-8}).
\end{align*}
A direct computation confirms $c_0+c_1+c_2+c_3=1$ and that each of the
three bracketed sums vanishes. Hence $A_N = L + O(N^{-8})$.
\end{proof}

\begin{corollary}[Four-Scale Polygon Acceleration]
\[
  \widetilde{P}_N^{(4)}
  = -\frac{1}{2835}P_N + \frac{4}{135}P_{2N}
    - \frac{64}{135}P_{4N} + \frac{4096}{2835}P_{8N}
  = \pi + O(N^{-8}).
\]
\end{corollary}

% ============================================================
\section{General $r$-Scale N-Series Acceleration}

\begin{theorem}[General Limit-Preserving Even-Power Acceleration]
\label{thm:general}
Let
\[
  a_N = L + \sum_{p=1}^{r}\frac{C_{2p}}{N^{2p}} + O(N^{-2r-2}),
\]
and let $m_0, m_1, \dots, m_r$ be distinct positive integers. Define
\[
  A_N = \sum_{j=0}^r c_j\,a_{m_jN},
  \qquad
  c_j = \prod_{\substack{i=0\\i\neq j}}^{r}
        \frac{m_j^2}{m_j^2 - m_i^2}.
\]
Then
\[
  A_N = L + O(N^{-2r-2}),\qquad
  \lim_{N\to\infty}A_N = L.
\]
\end{theorem}

\begin{proof}
For each $j$,
\[
  a_{m_jN} = L + \sum_{p=1}^{r}\frac{C_{2p}}{m_j^{2p}N^{2p}} + O(N^{-2r-2}).
\]
Hence
\[
  A_N = \Bigl(\sum_{j}c_j\Bigr)L
      + \sum_{p=1}^{r}\Bigl(\sum_{j}\frac{c_j}{m_j^{2p}}\Bigr)
        \frac{C_{2p}}{N^{2p}} + O(N^{-2r-2}).
\]
Setting $x_j = 1/m_j^2$, the coefficients $c_j$ are the Lagrange interpolation
weights at $x=0$ for the nodes $\{x_j\}$. They therefore satisfy
\[
  \sum_{j=0}^r c_j = 1,\qquad
  \sum_{j=0}^r c_j x_j^p = 0\quad\text{for }p=1,\dots,r.
\]
Hence all terms through $N^{-2r}$ cancel, leaving $A_N = L + O(N^{-2r-2})$.
The limit statement follows from $\sum c_j = 1$.
\end{proof}

\begin{corollary}[General $r$-Scale Polygon Acceleration]
For $P_N = N\sin(\pi/N)$ and distinct positive integers $m_0,\dots,m_r$,
\[
  \widetilde{P}_N^{(r)}
  = \sum_{j=0}^r c_j\,P_{m_jN}
  = \pi + O(N^{-2r-2}),\qquad
  \lim_{N\to\infty}\widetilde{P}_N^{(r)} = \pi.
\]
\end{corollary}

The convergence hierarchy is therefore:
\[
  P_N = \pi + O(N^{-2})
  \;\xrightarrow{r=1}\; O(N^{-4})
  \;\xrightarrow{r=2}\; O(N^{-6})
  \;\xrightarrow{r=3}\; O(N^{-8})
  \;\to\;\cdots
\]
with the exact limit $\pi$ preserved at every stage.

% ============================================================
\section{Geometric Interpretation of the N-Series Operator}

\subsection{Symmetry-Induced Even-Power Structure}

The even-power expansion of $P_N$ is not accidental. It arises from the
underlying symmetry of the chord construction: chord length is odd in the
subtended angle, but the chord-to-angle ratio is even. This symmetry
enforces the even-power error hierarchy and is the structural prerequisite
for N-series cancellation.

\subsection{Scaled Geometric Replication}

Each term $P_{m_jN}$ represents the same Archimedean construction at a
finer resolution. The family $\{P_{m_jN}\}$ is not a collection of
unrelated approximations but a hierarchy of self-similar geometric instances.

\subsection{Cancellation as Geometric Balancing}

The coefficients $c_j$ balance contributions from multiple scaled
constructions so that their leading deviations cancel. The operator
does not introduce new geometry; it redistributes existing geometric
information across scales.

\subsection{Interpolation Structure and Uniqueness}

Setting $x_j = 1/m_j^2$, the N-series coefficients are the Lagrange
weights at $x=0$ for the nodes $\{x_j\}$:
\[
  c_j = \prod_{\substack{i=0\\i\neq j}}^r \frac{x_i}{x_i - x_j}.
\]
Thus the N-series operator is \emph{uniquely determined} by the simultaneous
requirements of limit preservation and even-power cancellation.

\begin{remark}
Geometric constructions define approximation sequences, while symmetry determines
their error structure. The N-series operator reveals that these sequences admit
exact cancellation laws, allowing convergence to be accelerated without altering
the underlying constant.
\end{remark}

% ============================================================
\section{Symmetry Structure of the Trisection Refinement Law}

\subsection{Signed Angular Error}

Let $\theta$ be the target angle and $D$ the initial geometric approximation
to $\theta/3$. Define the signed error $\varepsilon = D - \theta/3$. Under
orientation reversal, $\theta \mapsto -\theta$ and $\varepsilon \mapsto -\varepsilon$,
so the error variable is odd.

\subsection{Symmetric Geometric Correction}

The refinement step constructs a corrected angle from a symmetric chord-based
configuration. Because the configuration is symmetric under orientation reversal,
the resulting correction function $\Phi(\varepsilon)$ satisfies
\[
  \Phi(-\varepsilon) = -\Phi(\varepsilon),
\]
so $\Phi$ is an odd function of the error.

\subsection{Cancellation of the Linear Term}

Expanding:
\[
  \Phi(\varepsilon) = a_1\varepsilon + a_3\varepsilon^3 + a_5\varepsilon^5 + \cdots.
\]
The reflection symmetry of the equal-length geometric components forces
$a_1 = 0$. With the linear term eliminated, the leading behavior is
\[
  \varepsilon_1 = a_3\varepsilon^3 + O(\varepsilon^5),
\]
and the refinement law takes the form
\[
  \varepsilon_{n+1} = -A\,\varepsilon_n^3 + O(\varepsilon_n^5),
  \qquad A > 0.
\]

\subsection{Unified Interpretation}

This mechanism parallels the N-series even-power structure:

\begin{center}
\begin{tabular}{lll}
\hline
\textbf{Setting} & \textbf{Mechanism} & \textbf{Result} \\
\hline
Polygon / $\pi$ & normalize odd chord law & even-power expansion \\
Trisection & symmetric correction of odd error & cubic odd refinement \\
\hline
\end{tabular}
\end{center}

Both arise from the same underlying principle: \emph{geometric symmetry
under orientation reversal determines the parity of the error algebra}.

\begin{remark}
The cubic refinement law in approximate angle trisection is a direct consequence
of symmetry-induced cancellation. It is not an artifact of approximation, but a
structural property of the underlying geometric construction.
\end{remark}

% ============================================================
\section{Conclusion}

We have established a complete framework for geometric acceleration of
convergent sequences with even-power error expansions. The key results are:

\begin{enumerate}
  \item The even-power structure of $P_N$ is forced by the odd symmetry of
        chord geometry under orientation reversal (Lemma 2.1, Proposition 2.2).
  \item The N-series accelerated limit equals $\pi$ exactly: $\pi_g = \pi$
        (Theorem 3.1).
  \item The acceleration operators are uniquely determined by the Lagrange
        interpolation structure at scaled nodes $x_j = 1/m_j^2$
        (Theorem 5.1).
  \item Each additional scale cancels one further even-power term, advancing
        convergence from $O(N^{-2r})$ to $O(N^{-2r-2})$ while preserving
        the exact limit.
  \item The same symmetry principle governs the cubic trisection refinement
        law, unifying the polygon and trisection branches of the framework.
\end{enumerate}

\begin{remark}
Geometry dictates the parity of error: normalization of an odd chord law
yields even-power asymptotics, while symmetric correction of signed angular
error yields higher-order odd refinement laws. These two mechanisms form the
structural foundation of the Baker N-series framework.
\end{remark}

\end{document}
